
\chapter{Prólogo}

\section{Objetivos de estos apuntes}

El curso estándar en un grado de matemáticas toma por dada las fundaciones que lo hacen posible, siendo $ ZFC $ la elección unánime. Estos apuntes sirven cómo intento de respuesta a la pregunta: ¿Es posible cursar el grado de matemáticas en fundaciones constructivas?.

Definamos pues que son las matemáticas constructivas:

\begin{quote}
Constructive mathematics is distinguished from its traditional counterpart, classical mathematics, by the strict interpretation of the phrase ``there exists'' as ``we can construct''. In order to work constructively, we need to re-interpret not only the existential 
quantifier but all the logical connectives and quantifiers as instructions on how to construct a proof of the statement involving these logical expressions. 

--- Stanford Encyclopedia of Philosophy
\end{quote}

Las matemáticas constructivas son entonces, utilizando la analogía entre topoi y universos matemáticos, matemáticas internas a topoi con buenas propiedades de computabilidad (o continuidad, para los geómetras). Es decir, se prohíbe el uso de cualquier principio 
sin contenido computacional.

Tómese la Ley del Tercero Excluido (LTE), la proposición que afirma $ \forall p \in Prop.\; p \vee \neg p $ . Computacionalmente, se nos pide ser capaces de realizar una elección, juzgar la veracidad de $ p $, por cada proposición existente. No es difícil ver por
qué dicha ley es completamente anti-computacional.

Bien es cierto que es demostrable que ciertos resultados son inalcanzables al constructivismo, por si renunciar a la Ley del Tercero Excluido no fuese poco, debemos asimismo renunciar al axioma de la elección pues este implica LTE. Es además un resultado muy conocido
que el axioma de la elección es equivalente a la proposición que afirma que todo espacio vectorial tiene una base. De hecho, muchas áreas de las matemáticas, si no todas, se complican si uno no es cuidadoso reformulando las definiciones y demostraciones clásicas.

La reacción natural ante tal posición es quizás de rechazo: ¿Por qué alguien querría limitarse a tal paronama sin motivo alguno? Debates sobre la filosofía de las matemáticas aparte, el desarrollo de las matemáticas constructivas en sí puede verse como desarrollo
de una área de interés para muchas ramas de la ciencia que sin embargo es renegada por muchos.

Asimismo, todas las matemáticas contenidas en estos apuntes son formalizadas en Cubical-Agda, extensión que añade ideas de Teoría Cúbica de Tipos a el asistente de demostraciones Agda. La Teoría Cúbica de Tipos da un modelo reciente donde el Axioma de Univalencia de
Voyevodski tiene una interpretación computacional.

\section{Matemáticas sin la Ley del Tercero Excluido}



\subsection{Computación}

\subsection{Limitaciones}
